\exposetitle{XIX}{Reductive groups: generalities}
\label{XIX}
\begin{center}\textit{by M.~Demazure}\end{center}

{\renewcommand{\thefootnote}{(0)}\nde{Version of 13/10/2024}}%
The remainder of this Seminar is devoted to the study of reductive groups. Its principal aim is the generalization of the classical results of Chevalley (\textit{Bible} and \textit{T\^ohoku}){\renewcommand{\thefootnote}{(1)}\nde{See the bibliographical references at the end of this Expos\'e. In particular, a new edition of the Chevalley Seminar 1956--58, cited [Bible], revised by P.~Cartier, was published in 2005, cf. [Ch05].}} to arbitrary base schemes, the two central results being the \emph{uniqueness theorem} (Exp.~XXIII, th.~4.1 and cor.~5.1 to 5.10) and the \emph{existence theorem} (Exp.~XXV, th.~1.1) for ``pinned'' reductive group schemes corresponding to prescribed ``root data''. The proofs used are inspired by those of Chevalley, the technique of schemes allowing them to be made more effective.

The results of the first volume of the \textit{Bible} (Exp.~1 to 13) will be used systematically. On the other hand, we shall prove directly over an arbitrary scheme the results of the second volume (in particular, the isomorphism theorem); knowledge of the proofs over an algebraically closed field is therefore not absolutely indispensable.

In the proof of these two fundamental results, we shall use only the most elementary results of the theory of groups of multiplicative type, contained for the most part in Exp.~VIII and IX; on the other hand, we shall make essential use of the results of Expos\'e~XVIII.{\renewcommand{\thefootnote}{(2)}\nde{More precisely, proposition XVIII 2.3 (extension of a ``generic homomorphism'' between groups) is used in XXII 4.1.11 and then in Exp.~XXIII (proof of the uniqueness theorem), and also in Exp.~XXIV; theorem XVIII 3.7 (construction of a group from a group germ) is used only in Exp.~XXV.}} The reader especially interested in the existence and uniqueness theorems may, on a first reading, skip Expos\'es~X to XVII.

\sgasection{1}{Reminders on groups over an algebraically closed field}
\label{XIX.sec.1}
\numpara{1.1} \emph{In this section, $k$ will always denote an algebraically closed field.} As announced above, the only results of \textit{Bible} used in what follows are in volume~1 (Expos\'es~1 to 13). All the results of \textit{Bible} (loc.~cit.) concerning \emph{semisimple groups} are valid more generally for \emph{reductive groups} (definition below), and their proof is identical, with the following harmless modifications:
\begin{itemize}
\item Exp.~9, \S~4, definition~3: see 1.6.1 below.
\item Exp.~12, \S~4, theorem~2, e): delete ``finite''.
\item Exp.~13, \S~3, theorem~2: replace ``rank'' by ``semisimple rank''.
\item Exp.~13, \S~4, corollary~2 to theorem~3: replace ``rank'' by ``reductive rank''.
\end{itemize}

\numpara{1.2} Let $G$ be a smooth affine connected $k$-group. The \emph{radical} of $G$ (\textit{Bible}, \S~9.4, prop.~2){\renewcommand{\thefootnote}{(3)}\nde{The reference 9-06 (= Exp.~9, p.~6) has been replaced by \S~9.4 (= Exp.~9, \S~4), which applies equally to [Bible] and [Ch05]. When the numbering of [Ch05] differs from [Bible], as is the case in Exp.~6, both references will be indicated explicitly.}} is the reduced subgroup associated with the identity component of the intersection of the Borel subgroups of $G$; it is also the largest smooth connected solvable invariant subgroup of $G$; we shall denote it by $\operatorname{rad}(G)$.

The \emph{unipotent radical} of $G$ is the unipotent part of the radical of $G$; it is also the largest smooth connected unipotent invariant subgroup of $G$; we shall denote it by $\operatorname{rad}^{u}(G)$.

\numpara{1.3} Let $G$ be a smooth affine connected $k$-group, and $Q$ a torus of $G$. Then the centralizer $\uCentr_G(Q)$ of $Q$ in $G$ is a closed subgroup of $G$ (Exp.~VIII, 6.7), smooth (Exp.~XI, 2.4) and connected (cf. \textit{Bible}, \S~6.6, th.~6 a) or [Ch05], \S~6.7, th.~7 a)).

One has the fundamental relation:
\begin{equation}\tag{1.3.1}\label{XIX.eq.1.3.1}
\operatorname{rad}^{u}(\uCentr_G(Q))=\operatorname{rad}^{u}(G)\cap\uCentr_G(Q).
\end{equation}
{\renewcommand{\thefootnote}{(4)}\nde{The following has been expanded.}}%
First, $U=\operatorname{rad}^{u}(G)\cap\uCentr_G(Q)$ is a unipotent invariant subgroup of $\uCentr_G(Q)$; let us show that it is smooth and connected. If one makes $Q$ act on $\operatorname{rad}^{u}(G)$ by inner automorphisms, $U$ is precisely the scheme of invariants of this action. Now this is smooth and connected, by lemma~1.4 below.

Consequently, $U$ is a closed subgroup of $\operatorname{rad}^{u}(\uCentr_G(Q))$. On the other hand, by \textit{Bible}, Exp.~12, \S~3, cor.~to th.~1, one has $\operatorname{rad}^{u}(\uCentr_G(Q))(k)\subset\operatorname{rad}^{u}(G)(k)$. Equality (1.3.1) follows.

Let us point out a particular case of the preceding result: if $G$ is a smooth affine connected $k$-group and $T$ is a maximal torus of $G$,
\begin{equation}\tag{1.3.2}\label{XIX.eq.1.3.2}
\uCentr_G(T)=T\cdot(\uCentr_G(T)\cap\operatorname{rad}^{u}(G)).
\end{equation}
Indeed (cf. \textit{Bible}, \S~6.6, th.~6 c) or [Ch05], \S~6.7, th.~7 c)), $\uCentr_G(T)$ is a solvable group, hence the semidirect product of its maximal torus $T$ and its unipotent radical.

By the density theorem (cf. \textit{Bible}, \S~6.5, th.~5 a) or [Ch05], \S~6.6, th.~6 a)), the union of the $\uCentr_G(T)$, as $T$ runs through the set of maximal tori of $G$, contains a dense open subset of $G$; it therefore follows from (1.3.2):
\begin{corollary}{1.3.3}\label{XIX.1.3.3}
If $G$ is a smooth affine connected $k$-group, the union of the $T\cdot\operatorname{rad}^{u}(G)$, where $T$ runs through the set of maximal tori of $G$, contains a dense open subset of $G$.
\end{corollary}

\begin{notation}{1.4.0}\label{XIX.1.4.0}
{\renewcommand{\thefootnote}{(5)}\nde{These notations, which appeared in 2.3 of the original, have been inserted here; on the other hand, the statement and proof of 1.4 have been expanded.}}%
We shall systematically use the following notation: if $S$ is a scheme and $s$ a point of $S$, $\overline{s}$ will denote the spectrum of an algebraic closure $\overline{\kappa(s)}$ of $\kappa(s)$.
\end{notation}

\begin{lemma}{1.4}\label{XIX.1.4}
{\renewcommand{\thefootnote}{(5)}\footnotemark}%
Let $S$ be a scheme, $Q$ an $S$-torus acting on a group $S$-scheme $H$, separated and smooth over $S$.
\begin{enumeratei}
\item The functor of invariants $H^Q$ is representable by a closed subscheme of $H$, smooth over $S$.
\item If in addition $H$ is affine over $S$, and has connected fibers, then so does $H^Q$.
\end{enumeratei}
\end{lemma}
First, by VIII 6.5 (d),{\renewcommand{\thefootnote}{(6)}\nde{See also the additions made in VI$_{\mathrm B}$, 6.2.4 (d) and 6.5.5 (a).}} since $H$ is separated over $S$ and $Q$ essentially free over $S$, $H^Q$ is representable by a closed subscheme of $H$. In particular, if $H$ is affine over $S$, then so is $H^Q$.

Consider now the semidirect product $G=H\cdot Q$; it is smooth and separated over $S$. Then the centralizer $\uCentr_G(Q)$ is representable by a closed subscheme of $G$, of finite presentation over $G$, by Exp.~XI 6.11 (a),{\renewcommand{\thefootnote}{(6)}\footnotemark} and it is smooth over $S$ by Exp.~XI 2.4.

Since $\uCentr_G(Q)=H^Q\cdot Q$, it follows that $H^Q$ is smooth over $S$. (Indeed, using the section $\sigma:H^Q\to\uCentr_G(Q)$ deduced from the unit section $\varepsilon_Q$ of $Q$, one sees at once that $H^Q$ is formally smooth over $S$; since moreover $\varepsilon_Q$ is of finite presentation, $H^Q$ is locally of finite presentation over $S$.) This proves (i).

Finally, suppose that $H$ is affine over $S$ and has connected fibers. Then each geometric fiber $G_{\overline{s}}$ of $G$ is a smooth affine connected $\overline{\kappa(s)}$-group, hence, by the first assertion of 1.3, the same is true of
\[
\uCentr_{G_{\overline{s}}}(Q_{\overline{s}})=(\uCentr_G(Q))_{\overline{s}}=(H^Q)_{\overline{s}}\cdot Q_{\overline{s}},
\]
and this implies that $(H^Q)_{\overline{s}}$ is connected.

\numpara{1.5} Recall that the \emph{reductive rank} of the smooth affine $k$-group $G$ is the common dimension of the maximal tori of $G$. We shall denote it by $\operatorname{rgred}(G/k)$ or $\operatorname{rgred}(G)$. For $\operatorname{rgred}(G/k)=0$, it is necessary and sufficient that $G$ be unipotent (i.e. that $G=\operatorname{rad}^{u}(G)$), by \textit{Bible}, \S~6.4, cor.~1 to th.~4 or [Ch05], \S~6.5, cor.~1 to th.~5.

If $H$ is an invariant subgroup of the smooth affine $k$-group $G$, the quotient $G/H$ is affine and smooth (Exp.~VI$_{\mathrm B}$, 11.17 and VI$_{\mathrm A}$, 3.3.2 (iii)). Moreover (\textit{Bible}, \S~7.3, th.~3, a) and c)), one has
\[
\operatorname{rgred}(G)=\operatorname{rgred}(G/H)+\operatorname{rgred}(H).
\]

\begin{definition}{1.6.1}\label{XIX.1.6.1}
{\renewcommand{\thefootnote}{(7)}\nde{Paragraph 1.6 has been changed into paragraphs 1.6.1 to 1.6.3. Note that in this Seminar every reductive (resp. semisimple, cf. 1.8) $k$-group is, by definition, connected.}}%
One says that the $k$-group $G$ is \emph{reductive} if it is smooth, affine and connected, and if $\operatorname{rad}(G)$ is a torus, that is, if $\operatorname{rad}^{u}(G)=\{e\}$.
\end{definition}
\begin{lemma}{1.6.2}\label{XIX.1.6.2}
Let $G$ be a reductive $k$-group.
\begin{enumeratei}
\item If $Q$ is a torus of $G$, then $\uCentr_G(Q)$ is reductive.
\item In particular, if $T$ is a maximal torus of $G$, then $\uCentr_G(T)=T$.
\item The center of $G$ is contained in every maximal torus, hence is diagonalizable.
\item The radical of $G$ is the (unique) maximal torus of $\uCentr(G)$.
\end{enumeratei}
\end{lemma}
Indeed, (i) follows from (1.3.1), (ii) from (1.3.2), and (iii) follows from (ii). Finally, the maximal torus of $\uCentr(G)$ (i.e. the identity component $\uCentr(G)^0$) is a smooth solvable connected subgroup, invariant in $G$, hence contained in $\operatorname{rad}(G)$; conversely, $\operatorname{rad}(G)$ is a torus invariant in $G$, hence central (\textit{Bible}, \S~4.3, cor.~to prop.~2), whence (iv).

\begin{remark}{1.6.3}\label{XIX.1.6.3}
If $G$ is reductive and $\dim(G)\ne\operatorname{rgred}(G)$, then $\dim(G)-\operatorname{rgred}(G)\geq2$. Indeed, this difference is always even (cf. 1.10 below).
\end{remark}

\numpara{1.7} Let $G$ be a smooth affine connected $k$-group and $H$ a smooth connected invariant subgroup. Then
\begin{equation}\tag{1.7.1}\label{XIX.eq.1.7.1}
\operatorname{rad}(H)=(\operatorname{rad}(G)\cap H)^0_{\mathrm{red}}
\quad\text{and}\quad
\operatorname{rad}^{u}(H)=(\operatorname{rad}^{u}(G)\cap H)^0_{\mathrm{red}},
\end{equation}
as one sees at once. In particular, if $G$ is reductive, so is $H$.

Let $f:G\to G'$ be a surjective morphism of smooth affine connected $k$-groups. Then
\begin{equation}\tag{1.7.2}\label{XIX.eq.1.7.2}
f(\operatorname{rad}^{u}(G))=\operatorname{rad}^{u}(G').
\end{equation}
In particular, if $G$ is reductive, so is $G'$.

Let us prove (1.7.2). First, $f$ sends $\operatorname{rad}^{u}(G)$ into $\operatorname{rad}^{u}(G')$. Introducing
\[H=(f^{-1}(\operatorname{rad}^{u}(G')))^0_{\mathrm{red}},\]
which contains $\operatorname{rad}^{u}(G)$, one has $\operatorname{rad}^{u}(H)=\operatorname{rad}^{u}(G)$ and one is reduced to the case where $G=H$, i.e. where $G'$ is unipotent. Since the union of the $T\cdot\operatorname{rad}^{u}(G)$ ($T$ running through the set of maximal tori of $G$) is dense in $G$, the union of the $f(T)f(\operatorname{rad}^{u}(G))$ is dense in $G'$; but $f(T)$ consists of semisimple elements, hence $f(T)=\{e\}$, since $G'$ is unipotent; this shows that $f(\operatorname{rad}^{u}(G))$ is dense in $G'$. Therefore, by \textit{Bible}, \S~5.4, lemma~4 or [Ch05], \S~6.1, lemma~1,{\renewcommand{\thefootnote}{(8)}\nde{The following has been expanded.}} $f(\operatorname{rad}^{u}(G))$ is an open subgroup of $G'$; since the latter is connected, it follows that $f(\operatorname{rad}^{u}(G))=G'$.

(N.~B.: under the same hypotheses one can prove that $f(\operatorname{rad}(G))=\operatorname{rad}(G')$.)

\numpara{1.8} One says that the $k$-group $G$ is \emph{semisimple} if it is smooth, affine and connected and if $\operatorname{rad}(G)=\{e\}$, that is, if $G$ is reductive and $\uCentr(G)$ finite. If $G$ is any smooth affine connected $k$-group, then $G/\operatorname{rad}(G)$ is semisimple (\textit{Bible}, \S~9.4, prop.~2), and $G/\operatorname{rad}^{u}(G)$ is reductive. One calls \emph{semisimple rank} of $G$, and denotes by $\operatorname{rgss}(G/k)$ or $\operatorname{rgss}(G)$, the reductive rank of $G/\operatorname{rad}(G)$.

If $G$ is reductive, one therefore has
\[
\operatorname{rgred}(G)=\operatorname{rgss}(G)+\dim(\operatorname{rad}(G)).
\]
If $G$ is a smooth affine connected $k$-group and $Q$ is a central subtorus of $G$, then $G/Q$ is semisimple if and only if $G$ is reductive and $Q=\operatorname{rad}(G)$. Indeed, if $G/Q$ is semisimple, $Q\supset\operatorname{rad}(G)$, hence $\operatorname{rad}(G)$ is a torus, hence $G$ is reductive, hence $\operatorname{rad}(G)$ is the maximal torus of $\uCentr(G)$, hence $\operatorname{rad}(G)=Q$. If $G$ is reductive and $Q$ is a central subgroup, then ($G/Q$ is reductive and) $\operatorname{rgss}(G)=\operatorname{rgss}(G/Q)$.

\numpara{1.9} If $K$ is an algebraically closed extension of $k$ and $G$ is a smooth affine connected $k$-group, $G$ is reductive (resp. semisimple) if and only if $G_K$ is, and one has
\[
\operatorname{rgred}(G/k)=\operatorname{rgred}(G_K/K),\qquad
\operatorname{rgss}(G/k)=\operatorname{rgss}(G_K/K).
\]

\numpara{1.10} Let $G$ be a smooth connected $k$-group and $T$ a torus of $G$.{\renewcommand{\thefootnote}{(9)}\nde{The following sentence has been added.}} One denotes by $\mathfrak g$ the Lie $k$-algebra of $G$, i.e. $\mathfrak g=\Lie(G)=\Lie(G/k)(k)$; similarly, one puts $\mathfrak t=\Lie(T)$. Then $\mathfrak g$ decomposes under the action of $T$ (by $T\to G$ and the adjoint action of $G$) into
\begin{equation}\tag{1.10.1}\label{XIX.eq.1.10.1}
\mathfrak g=\mathfrak g^0\oplus\coprod_{\alpha\in R}\mathfrak g^\alpha,
\end{equation}
where the $\alpha\in R$ are nontrivial characters of $T$ and the $\mathfrak g^\alpha$ are $\ne0$. The characters $\alpha\in R$ are called the \emph{roots} of $G$ with respect to $T$. By Exp.~II, 5.2.3, (i), one has
\begin{equation}\tag{1.10.2}\label{XIX.eq.1.10.2}
\mathfrak g^0=\Lie(\uCentr_G(T)).
\end{equation}
In particular,{\renewcommand{\thefootnote}{(10)}\nde{The following has been expanded.}} since $\uCentr_G(T)$ is connected (1.3 when $G$ is affine, and Exp.~XII 6.6 b) in the general case), $T$ is its own centralizer if and only if $\mathfrak g^0=\Lie(T)$.

This condition implies that $T$ is maximal and that $\uCentr(G)\subset T$, hence by Exp.~XII 8.8 d) one has:
\begin{equation}\tag{1.10.3}\label{XIX.eq.1.10.3}
\uCentr(G)=\Ker(T\to\GL(\mathfrak g))=\bigcap_{\alpha\in R}\Ker(\alpha);
\end{equation}
moreover, $\uCentr(G)$ is then affine, hence $G\to G/\uCentr(G)$ affine; since the latter group is affine (Exp.~XII 6.1),{\renewcommand{\thefootnote}{(11)}\nde{See also the addition VI$_{\mathrm B}$, 6.2.6.}} so is $G$.

When $G$ is reductive and $T$ maximal, the roots in the preceding sense coincide with the roots in the sense of \textit{Bible}, \S~12.2, def.~1; the latter are indeed roots in the sense of this paragraph (\textit{Bible}, \S~13.2, th.~1, c)) and there are $\dim(G)-\dim(T)$ of them (\textit{Bible}, \S~13.4, cor.~2 to th.~3). Moreover, if $\alpha$ is a root, so is $-\alpha$ (\textit{Bible}, \S~12.2, cor.~to prop.~1). (As usual, one writes the group structure of $\Hom_{k\text{-gr.}}(T,\mathbf{G}_{\mathrm m,k})$ indifferently additively or multiplicatively.) It follows that, for $G$ reductive,
\begin{equation}\tag{1.10.4}\label{XIX.eq.1.10.4}
\dim(G)-\operatorname{rgred}(G)=\operatorname{Card}(R)
\end{equation}
is always even.

The semisimple rank of the reductive group $G$ is the rank of the subset $R$ of the $\QQ$-vector space $\Hom_{k\text{-gr.}}(T,\mathbf{G}_{\mathrm m,k})\otimes\QQ$ (\textit{Bible}, \S~13.3, th.~2).

\begin{lemma}{1.11}\label{XIX.1.11}
Let $k$ be an algebraically closed field, $G$ a smooth affine connected $k$-group, $T$ a torus of $G$, $W(T)=\uNorm_G(T)/\uCentr_G(T)$ the Weyl group of $G$ with respect to $T$. The following conditions are equivalent:
\begin{enumeratei}
\item $G$ is reductive, $T$ is maximal, $\operatorname{rgss}(G)=1$.
\item $G$ is reductive, $T$ is maximal, $G\ne T$; there exists a subtorus $Q$ of $T$, of codimension~1 in $T$, central in $G$.
\item $G$ is not solvable and $\dim(G)-\dim(T)\leq2$.
\item $W(T)\ne\{e\}$ and $\dim(G)-\dim(T)\leq2$.
\item $W(T)=(\ZZ/2\ZZ)_k$ and $\dim(G)-\dim(T)=2$.
\end{enumeratei}
Moreover, under these conditions, there are exactly two roots of $G$ with respect to $T$; they are opposite. Under the conditions of (ii), $Q=\operatorname{rad}(G)$.
\end{lemma}
One clearly has (v) $\Rightarrow$ (iv). One has (iv) $\Rightarrow$ (iii) by \textit{Bible}, \S~6.1, cor.~3 to th.~1 or [Ch05], \S~6.2, cor.~3 to th.~2. Let us prove (iii) $\Rightarrow$ (ii): let $U$ be the unipotent radical of $G$; one knows that $G/U$ is reductive and is not a torus (otherwise $G$ would be solvable); hence, by (1.10.4), one has
\[
\dim(G/U)-\operatorname{rgred}(G/U)=\operatorname{Card}(R)\geq2;
\]
but
\[
\operatorname{rgred}(G)=\operatorname{rgred}(G/U)+\operatorname{rgred}(U)=\operatorname{rgred}(G/U),
\]
whence
\[
\begin{split}
\dim(G)-\dim(U)-\operatorname{rgred}(G)&=\dim(G/U)-\operatorname{rgred}(G/U)\\
&\geq2\geq\dim(G)-\dim(T)\geq\dim(G)-\operatorname{rgred}(G).
\end{split}
\]
This gives $\dim(U)=0$, hence $G$ is reductive, $\operatorname{rgred}(G)=\dim(T)$, hence $T$ is maximal, $\dim(G)-\dim(T)=2$. By \textit{Bible}, \S~10.4, prop.~8, there exists a singular torus $Q$ of codimension~1 in $T$; then $\uCentr_G(Q)$ is reductive (1.6.2 (i)), not solvable (definition of a singular torus), hence $\dim(\uCentr_G(Q))-\operatorname{rgred}(G)\geq2$, which proves that $G=\uCentr_G(Q)$ and $Q$ is central in $G$.

Let us prove (ii) $\Rightarrow$ (i). One knows that $G/Q$ is reductive (1.7) and that $\operatorname{rgred}(G/Q)=1$, hence $\operatorname{rgss}(G/Q)=0$ or $1$. The first case is impossible, for it would imply $\operatorname{rgss}(G)=0$, hence $G=T$; one therefore has $\operatorname{rgred}(G/Q)=\operatorname{rgss}(G/Q)=1$, which proves that $G/Q$ is semisimple; hence $Q$ is the radical of $G$ and $\operatorname{rgss}(G)=1$.
